\chapter{Implementation}

This chapter describes how each component was built. The source is organised as a
Cargo workspace with three Rust crates --- \texttt{signaling-server},
\texttt{native-host} and \texttt{host-service} --- alongside Tauri applications
for the attended host and the controller. The workspace deliberately excludes the
Tauri applications from its shared build profile so that their front-end
toolchains do not conflict with the systems crates.

\section{Project Layout}

\begin{table}[H]
\centering
\caption{Repository structure (principal paths)}
\label{tab:layout}
\small
\begin{tabular}{@{}p{5cm}p{8.4cm}@{}}
\toprule
\textbf{Path} & \textbf{Contents} \\
\midrule
\texttt{crates/signaling-server} & The authentication and rendezvous server. \\
\texttt{crates/native-host}      & The headless capture, encode, stream and input host. \\
\texttt{crates/host-service}     & The Windows service that supervises the host in the active session. \\
\texttt{apps/host-agent}         & Attended Tauri host (WebView2, \texttt{getDisplayMedia}). \\
\texttt{apps/controller-app}     & Tauri controller application. \\
\texttt{src/}, \texttt{shared/}  & Signaling logic and the shared client used during early milestones. \\
\texttt{vendor/webrtc-dtls}      & Patched DTLS crate (see Section~5.6). \\
\bottomrule
\end{tabular}
\end{table}

\section{Signaling Server}

The server is an asynchronous Tokio program that accepts WebSocket connections
through \texttt{tokio-tungstenite}. Each connection is driven by a small state
machine: it is created in an unauthenticated state, a challenge is issued
immediately, and no message other than \texttt{auth\_response} is processed until
a valid signature has been verified with \texttt{ed25519-dalek}. Verified
connections are recorded in an in-memory registry keyed by public key and, for
hosts, by connect identifier, so that a controller's \texttt{connect\_request}
can be routed to the correct socket.

TLS is optional and, when enabled, is terminated in-process by
\texttt{tokio-rustls} using the \texttt{ring} provider; the server reads a
PKCS\#12 identity whose path is supplied through the environment. Persistence is
likewise optional: when a database URL is configured, enrolled public keys and a
per-session audit record are written to PostgreSQL through \texttt{sqlx}. Neither
feature is required to run the system, which keeps the development loop light.

\begin{lstlisting}[language=Rust,caption={Authentication gate, in essence.}]
// A connection is unauthenticated until it proves possession of an enrolled key.
let nonce = random_nonce();
send(&mut ws, json!({ "type": "auth_init", "nonce": b64(&nonce) })).await?;

let resp = recv(&mut ws).await?;              // expect auth_response
let pubkey  = decode_pubkey(&resp["pubkey"])?;
let sig     = decode_sig(&resp["sig"])?;

if !enrolled.contains(&pubkey) { return reject("unknown key"); }
pubkey.verify_strict(&nonce, &sig)            // reject on any error
      .map_err(|_| reject("bad signature"))?;
// ...only now is the socket promoted to authenticated and allowed to register.
\end{lstlisting}

\section{Native Host --- Capture and Encode}

The native host is the component that removes the browser from the loop.
Framebuffer acquisition uses the \texttt{scrap} crate, which wraps the Desktop
Duplication API; each captured frame arrives as a BGRA buffer with a stride that
must be respected when repacking. Frames are converted to the encoder's expected
planar format and passed to OpenH264 through the \texttt{openh264} crate, which
emits H.264 access units. This capture happens with no window, no permission
prompt and no source picker --- the defining property recorded as FR-8 --- because
a native process reads the composited desktop surface directly rather than
requesting it through \texttt{getDisplayMedia}.

The first milestone of the native host wrote a single captured frame to
\texttt{native-capture.png} and a short encoded sequence to
\texttt{native-capture.h264} to verify the capture and encode stages in
isolation before any network code existed. This staged validation, one subsystem
proven before the next was attempted, is reflected throughout the commit
history.

\section{Native Host --- WebRTC and Signaling}

Encoded access units are written onto a video track created through
\texttt{webrtc-rs}. The host builds a peer connection, adds the track, and drives
the same signaling protocol as the graphical host: it authenticates with its
Ed25519 key, registers under its connect identifier, and on an accepted request
generates the SDP \texttt{offer}, consumes the controller's \texttt{answer}, and
trickles ICE candidates. Because the host runs headless, acceptance of an
unattended session is automatic once the presented password matches the
configured unattended secret; there is no dialogue to click.

\begin{lstlisting}[language=Rust,caption={Registering the native host after authentication.}]
send(&ws, json!({
    "type": "register",
    "connect_id": self.connect_id,     // nine-digit identifier
    "role": "host",
})).await?;
// then: await connect_request -> session_response(accept)
//       -> create offer -> set_local_description -> relay via server
\end{lstlisting}

\section{Native Host --- Input Injection}

Control is the reverse path. Pointer and keyboard events sent by the controller
arrive as messages that the host translates into synthetic operating-system
input. Two mechanisms are present: the cross-platform \texttt{enigo} crate for
ordinary desktops, and a direct Win32 \texttt{SendInput} path for the cases where
finer control or secure-desktop reach is required. Coordinates are mapped from
the controller's view geometry to the host's display resolution before
injection, so that a click lands where the operator intended regardless of any
scaling applied to the video.

\section{The DTLS Curve-Negotiation Patch}

A concrete obstacle deserves recording because it shaped the workspace. Against
current versions of Chrome, the DTLS handshake failed with an ``invalid named
curve'' error. The cause was that recent Chrome offers a post-quantum key-exchange
group first in its list, and the version of \texttt{webrtc-dtls} in use attempted
to honour the browser's first preference rather than selecting the first group
that both sides actually support. The fix, carried as a vendored patch in
\texttt{vendor/webrtc-dtls} and wired in through a Cargo \texttt{[patch]} section,
changes the negotiation to pick the first \emph{mutually} supported ECDHE curve.
This is the kind of defect that closed transports hide from their users and open
ones expose to repair; it is cited in the acknowledgements for that reason.

\begin{lstlisting}[language=Rust,caption={Workspace patch redirecting webrtc-dtls to the fixed local copy.}]
[patch.crates-io]
webrtc-dtls = { path = "vendor/webrtc-dtls" }
\end{lstlisting}

\section{Windows Service and Session Handling}

Unattended, pre-login presence is realised by registering the host as a Windows
service. The \texttt{host-service} crate uses the \texttt{windows-service} and
\texttt{windows} crates to install, start, stop and uninstall itself, and to run
as SYSTEM in session~0. From that vantage it detects the active console session
and ensures a capture process is running in whichever session and desktop is
currently active, relaunching it across logon, lock and session-switch
transitions. This is the mechanism by which the host is available from boot and
can, in principle, reach the secure desktop on which the login screen and UAC
prompts are drawn --- surfaces that no user-session application can capture.

Installation is a single command (\texttt{host-service.exe install} followed by
\texttt{sc start UpDeskHost}), after which the host survives reboot without
further intervention, satisfying FR-10.

\section{Administrative Tooling}

A small administrative command lists enrolled devices and identities with their
nine-digit connect identifiers, so that an operator can see which machines are
registered and address them. The native host additionally supports an
\texttt{id} command that prints its own connect identifier and unattended
password, which is the information a remote operator needs in order to connect.
